\newcommand{\eqPhaseRule}{
    \begin{equation}
        \label{eq:PhaseRule}
        \freedomDegree = \numberOfComponent - \numberOfPhases + 2
    \end{equation}
}


\newcommand{\eqDoubleEosVaporGas}{
  \begin{subequations}
    \label{eq:DoubleEosVaporGas} % Etiqueta para el grupo completo (opcional)
    \begin{align}
        {\pVapor}{\volume[]{}[]} &= {\zVapor}{\vaporMole}{\universalConstantGas}{\tVapor}  \quad \text{(Para el Vapor)} \label{eq:eosVapor} \\
        {\pGas}{\volume[]{}[]}   &= {\zGas}{\gasMole}{\universalConstantGas}{\tGas}  \quad \text{(Para el Gas)} \label{eq:eosGas}
    \end{align}
  \end{subequations}
}


\newcommand{\eqnClausiusClapeyron}{
    \begin{equation}
        \label{eqn:ClausiusClapeyron}
        \frac{d \pressure[]{}}{d \T} = \frac{L}{\T \Delta \volume[]{}}
    \end{equation}%
    Donde \absoluteHumidity es la humedad absoluta, y se representa a la humedad absoluta saturada como \saturatedAbsoluteHumidity, el cual es la máxima cantidad de vapor que el aire puede contener a la misma presión y temperatura.

}


% Con número: \eqnTotalHeat
% Sin número: \eqnTotalHeat*
\NewDocumentCommand{\eqnTotalHeat}{s}{
    \begin{equation}
        \IfBooleanTF{#1}{\nonumber}{\label{eq:heatFlow}}
        \totalHeat{} = \sensibleHeat{} + \latentHeat{}
    \end{equation}
}

\newcommand{\eqnHeatFlow}{
    \begin{equation}
        \label{eqn:HeatFlow}
        \heatFlow = {\specificHeatAtP}{\Delta \T}
    \end{equation}%
    Donde 

}



